% This is where all the commands should go that you want to define yourself.


% extra packages

% %%%%%%%%%%%%%%%%%%%%%%%%%%%%%%%%%%%%%%%%%%%%%%%%%%%%%%%%%%%%%%%%%%%%%%%%%%%%%%%%%%%%%%%%%%%%
 % For professional tables
\usepackage{dcolumn}% Align table columns on decimal point
\usepackage{booktabs} % if you're using \toprule, \midrule, \bottomrule
\usepackage{tabularx}
\usepackage{multirow}
\renewcommand{\arraystretch}{1.5}
\usepackage{threeparttable}
% %%%%%%%%%%%%%%%%%%%%%%%%%%%%%%%%%%%%%%%%%%%%%%%%%%%%%%%%%%%%%%%%%%%%%%%%%%%%%%%%%%%%%%%%%%%%

% %%%%%%%%%%%%%%%%%%%%%%%%%%%%%%%%%%%%%%%%%%%%%%%%%%%%%%%%%%%%%%%%%%%%%%%%%%%%%%%%%%%%%%%%%%%%
% for chemical forumals
\usepackage{chemformula}
% %%%%%%%%%%%%%%%%%%%%%%%%%%%%%%%%%%%%%%%%%%%%%%%%%%%%%%%%%%%%%%%%%%%%%%%%%%%%%%%%%%%%%%%%%%%%


\usepackage{gensymb} % for degree ° symbol
% \usepackage{mycommands}
% \usepackage{float}
%\usepackage{floatrow}
% \usepackage{dblfloatfix}
% \usepackage{placeins}

% %%%%%%%%%%%%%%%%%%%%%%%%%%%%%%%%%%%%%%%%%%%%%%%%%%%%%%%%%%%%%%%%%%%%%%%%%%%%%%%%%%%%%%%%%%%%
% SI unit
\usepackage{siunitx}
\DeclareSIUnit\torr{torr}
\DeclareSIUnit\ppm{ppm}
\DeclareSIUnit\ppb{ppb}
\DeclareSIUnit\bar{bar}
\DeclareSIUnit\gauss{G}
\DeclareSIUnit\neV{\nano\electronvolt}
\DeclareSIUnit\peV{\pico\electronvolt}
\DeclareSIUnit\kms{km/s}
\sisetup{print-unity-mantissa=false}  % To display only the exponent (i.e., 10⁻¹¹) 
\sisetup{separate-uncertainty=false} % use parenthesis style for uncertainty
% \sisetup{uncertainty-mode=separate} % use +/- style for uncertainty
\sisetup{uncertainty-mode=compact-marker} % turn this on to use style like 123.45(1.20)
% otherwise we  turn this on to use style like 123.45(120)
% %%%%%%%%%%%%%%%%%%%%%%%%%%%%%%%%%%%%%%%%%%%%%%%%%%%%%%%%%%%%%%%%%%%%%%%%%%%%%%%%%%%%%%%%%%%%

% \usepackage[normalem]{ulem}
\usepackage{changes}